\section*{Bab \ref{ch:g5:weather-and-climate} --- Cuaca dan Iklim}

\begin{solution}{exo:g5:weather-and-climate:1}
(a) \omterm{def:g1:watching-the-weather:weather}{Cuaca} (hanya pagi ini). (b) \omterm{def:g5:weather-and-climate:climate}{Iklim} (setiap musim panas). (c) \omterm{def:g1:watching-the-weather:weather}{Cuaca}
(kemarin sore). (d) \omterm{def:g5:weather-and-climate:climate}{Iklim} (setiap bulan, dari tahun ke tahun).
\end{solution}

\begin{solution}{exo:g5:weather-and-climate:2}
\omterm{def:g5:weather-and-climate:climate}{Iklim} suatu tempat adalah \omterm{def:g1:watching-the-weather:weather}{cuaca} yang biasanya terjadi di sana, bulan demi
bulan, selama bertahun-tahun. Para ahli \omterm{def:g1:watching-the-weather:weather}{cuaca} merata-ratakan hasil
pengukuran selama 30 tahun.
\end{solution}

\begin{solution}{exo:g5:weather-and-climate:3}
Bulan yang paling dingin adalah Desember (kira-kira 6~$^\circ$C; Januari
hampir sama dinginnya). Bulan yang paling basah adalah November (kira-kira
160~mm). Tidak: kurvanya tidak pernah turun di bawah 0~$^\circ$C.
\end{solution}

\begin{solution}{exo:g5:weather-and-climate:4}
Juli adalah bulan yang paling hangat, kira-kira 5~$^\circ$C. Tidak ada
bulan yang suhunya lebih dari 10~$^\circ$C.
\end{solution}

\begin{solution}{exo:g5:weather-and-climate:5}
\omterm{def:g3:continents-and-oceans:equator}{Khatulistiwa}: zona panas. \omterm{def:g3:continents-and-oceans:poles}{Kutub Utara}: zona dingin (yang di utara). Di
pertengahan antara \omterm{def:g3:continents-and-oceans:equator}{khatulistiwa} dan \omterm{def:g3:continents-and-oceans:poles}{Kutub Selatan}: zona sedang selatan
(zona itu terletak di antara \omterm{def:g5:weather-and-climate:tropics}{Garis Balik} Selatan dan Lingkaran Antarktik).
\end{solution}

\begin{solution}{exo:g5:weather-and-climate:6}
Hilo: zona panas. Phoenix dan Seattle: zona sedang utara. Utqiaġvik: zona
dingin utara, di luar Lingkaran Arktik.
\end{solution}

\begin{solution}{exo:g5:weather-and-climate:7}
Satu tahun penuh: $147 + 96 + 106 + 81 + 48 + 37 + 15 + 25 + 41 + 99 + 160 + 145
= 1\,000$ milimeter. Desember, Januari, dan Februari:
$145 + 147 + 96 = 388$ milimeter, lebih dari sepertiga \omterm{def:g4:clouds-rain-snow:precipitation}{hujan} setahun dalam
tiga bulan saja.
\end{solution}

\begin{solution}{exo:g5:weather-and-climate:8}
Phoenix: $35 - 14 = 21$ derajat. Hilo: $25 - 22 = 3$ derajat. Musim jauh
lebih banyak mengubah suhu di Phoenix; di Hilo suhunya hampir sama
sepanjang tahun.
\end{solution}

\begin{solution}{exo:g5:weather-and-climate:9}
\omterm{def:g5:weather-and-climate:climate}{Iklim} \omterm{def:g5:weather-and-climate:rainforest}{hutan hujan}: udaranya panas sepanjang tahun (27 atau 28~$^\circ$C,
tanpa musim dingin), dan banyak \omterm{def:g4:clouds-rain-snow:precipitation}{hujan} turun setiap bulan (paling sedikit
160~mm, dan lebih dari 2\,000~mm dalam setahun). Kota itu paling mungkin
terletak di zona panas, di dekat \omterm{def:g3:continents-and-oceans:equator}{khatulistiwa}.
\end{solution}

\begin{solution}{exo:g5:weather-and-climate:10}
Penjelasan pertama: pantai itu terletak di zona panas, yang hangat
sepanjang tahun, sedangkan kebun yang bersalju terletak di zona sedang atau
zona dingin, tempat bulan Januari jatuh pada musim dingin. Penjelasan
kedua: pantai itu terletak di \omterm{def:g3:continents-and-oceans:equator}{belahan bumi selatan}, tempat bulan Januari
jatuh pada musim panas, sedangkan kebun yang bersalju terletak di \omterm{def:g3:continents-and-oceans:equator}{belahan
bumi utara}, tempat bulan Januari jatuh pada musim dingin: musim-musimnya
terbalik di antara kedua belahan \omterm{def:g3:continents-and-oceans:globe}{bola dunia}.
\end{solution}

\begin{solution}{exo:g5:weather-and-climate:11}
Suhu: Phoenix jauh lebih panas, kira-kira 35~$^\circ$C pada bulan Juli,
dibandingkan dengan kira-kira 20~$^\circ$C pada bulan Agustus di Seattle,
dan bulan Januarinya (kira-kira 14~$^\circ$C) sudah sehangat bulan Mei di
Seattle. \omterm{def:g4:clouds-rain-snow:precipitation}{Hujan}: kurang dari 200~mm setahun di Phoenix, sebuah \omterm{def:g5:weather-and-climate:desert}{gurun},
dibandingkan dengan kira-kira 1\,000~mm di Seattle. Jadi \omterm{def:g5:weather-and-climate:climate}{iklim} keduanya
tidak sama. \omterm{def:g5:weather-and-climate:climate-zone}{Zona-zona iklim} memberi tahu kita, secara kasar, seberapa
banyak Matahari menghangatkan suatu tempat, dari \omterm{def:g3:continents-and-oceans:equator}{khatulistiwa} sampai ke
\omterm{def:g3:continents-and-oceans:poles}{kutub}; zona-zona itu tidak memberi tahu kita seberapa banyak \omterm{def:g4:clouds-rain-snow:precipitation}{hujan} yang
turun, atau pengaruh \omterm{def:g2:water-on-earth:ocean}{laut}, \omterm{def:g4:world-physical-map:range}{pegunungan}, dan angin.
\end{solution}
